\section{Partecipazione a progetti nazionali}
%
\outerlist{

\entrybig
	{\textbf{ECOSISTER}}{AimageLab, UNIMORE}
	{\textcolor{darkgray}{Ecosistema territoriale di innovazione dell’Emilia-Romagna} [\href{https://ecosister.it/spoke/spoke-4/}{link}]}{2023 -- in corso}
\innerlist{
	\entry{\textbf{Attività:} partecipazione alle attività di ricerca dello Spoke 4, con focus su sistemi di \textit{Artificial Intelligence} e \textit{Computer Vision} per applicazioni di \textit{smart city}, mobilità intelligente e analisi dei flussi pedonali.}
    \entry{\textbf{Finanziamento:} Progetto PNRR, finanziato dall’Unione Europea – \textit{NextGenerationEU}, nell’ambito del Piano Nazionale di Ripresa e Resilienza, tramite il Ministero dell’Università e della Ricerca (MUR).}

}

\entrybig
	{\textbf{LEGO.AI}}{AimageLab, UNIMORE}
	{\textcolor{darkgray}{LEarning the Geometry of knOwledge in AI systems} [\href{https://aimagelab.ing.unimore.it/imagelab/project.asp?idprogetto=99}{link}]}{2022 -- 2025}
\innerlist{
	\entry{\textbf{Attività:} studio e sviluppo di modelli di \textit{Geometric Deep Learning}, con particolare attenzione all’introduzione di vincoli e regolarizzatori geometrici per \textit{Continual Learning}, \textit{Few-Shot Learning} e \textit{Domain Adaptation}. Progettazione di modelli con maggiore capacità di generalizzazione e ridotto fabbisogno di dati annotati.}
    \entry{\textbf{Finanziamento:} Progetto PRIN (Progetto di Ricerca di Interesse Nazionale), finanziato dal Ministero dell’Istruzione, dell’Università e della Ricerca (MIUR).}
}

\entrybig
	{\textbf{AI 4 Vector}}{AimageLab, UNIMORE}
	{\textcolor{darkgray}{Artificial Intelligence and Remote Sensing}}{2019 -- 2021}
\innerlist{
    \entry{\textbf{Attività:} sviluppo di modelli predittivi per il monitoraggio e la previsione della proliferazione di insetti vettori di patologie infettive (\textit{Vector-Borne Diseases}), mediante l’analisi di variabili climatiche e ambientali derivate da immagini satellitari multi-banda ad alta risoluzione tramite \textit{Convolutional Neural Networks}.}
    \entry{\textbf{Finanziamento:} Ministero Italiano della Salute, progetto di ricerca corrente \textit{“Artificial Intelligence and Remote Sensing: innovative methods for monitoring the vectors and the associated ecological/environmental variables”}.}
}


\entrybig
	{\textbf{UBINVB}}{AimageLab, UNIMORE}
	{\textcolor{darkgray}{Ubiquitous objective measures of intergroup nonverbal behaviors}}{2018}
\innerlist{
	\entry{\textbf{Attività:} analisi automatica del comportamento umano e individuazione di costrutti psicologici emergenti in interazioni diadiche mediante l’impiego di \textit{Graph Neural Networks}, includendo la progettazione e lo sviluppo di modelli neurali basati su grafi per l’analisi di dati strutturati e relazionali.}
	\entry{\textbf{Finanziamento:} progetto “\href{http://www.cittaeducante.it/SitePages/sito/ceindex.aspx}{La Città Educante}”, Cluster Tecnologico Nazionale Smart Communities Tech, co-finanziato dal MIUR.}
}

}
